\documentclass[a4paper]{article}
\usepackage{amsfonts}
\usepackage{amsmath}
\usepackage{amssymb}
\usepackage{graphicx}     

\begin{document}
  
\noindent \large Auteur: Sabawoon Enayat \\
\noindent \large Studierichting: Informatica\\
\noindent \large Oefeningenreeks 6A nr. 7.5 \\

\medskip

\normalsize

Opgave: $x_1 \cdot x_3 = 36$ en $x^2 + x^4 = 60$

$\Rightarrow \left\{
	\begin{array}{lcr}
		x_1 \cdot (x_1 \cdot q^2)   & = & 36\\
		x_1 \cdot q + x_1 \cdot q^3 & = & 36\\
	\end{array}
\right. \Rightarrow \left\{
	\begin{array}{lcr}
		(x_1 \cdot q)^2 & = & 36\\
		x_1 (q + q^3)   & = & 60\\
	\end{array}
\right.$\\

$\Rightarrow \left\{
	\begin{array}{lcr}
		x_1 \cdot q              & = & 6\\
		x_1 \left(q + q^3\right) & = & 60\\
	\end{array}
\right. \Rightarrow \left\{
	\begin{array}{lcr}
		x_1                            & = & \dfrac{6}{q}\\
		x_1 \cdot \left(q + q^3\right) & = & 60 \\
	\end{array}
\right.$\\

$\Rightarrow \dfrac{6}{q} \left(q + q^3\right) = 60 \Rightarrow 6q^2 + 6 = 60 \Rightarrow q^2 + 1 = 10 \Rightarrow q^2 = 9$\\

$\Rightarrow q = \pm 3 \Rightarrow x_1 = \dfrac{6}{q} = 2$ of $x_1 = -2$\\

$\Rightarrow x_2 = x_1 \cdot q = 6$\\

$\Rightarrow x_3 = x_2 \cdot q = 18$ of $x_3 = x_2 \cdot q = -18$\\

$\Rightarrow (2, 6, 18, ...)$ of $(-2, 6, -18, ...)$\\

\begin{thebibliography}{999}
%\bibitem{a} Referentie
\end{thebibliography}
\end{document} 

